\documentclass[10pt,a4paper]{amsart}
\usepackage[margin=19mm]{geometry}
\usepackage{amsmath,amssymb,mathtools}
\usepackage[scr=rsfs,cal=euler]{mathalfa}
\usepackage{xfrac}
\usepackage[unicode,hidelinks,bookmarks=false]{hyperref}
\newcommand{\ie}{i.e.\ }
\newcommand{\cf}{cf.\ }
\newcommand{\resp}{respectively\ }
\newcommand{\Subsection}{\subsection}
\providecommand{\nobreakdash}{\nobreak\hbox{-}}
\setlength{\emergencystretch}{2em}
\multlinegap0pt
\usepackage{amssymb}
\usepackage{csquotes}
\usepackage{mathtools}
\usepackage[shortlabels]{enumitem}
\usepackage{xcolor}
\usepackage{tikz}
\newtheorem{lemma}{Lemma}
\newcommand{\N}{\mathbb{N}}
\newcommand{\R}{\mathbb{R}}
\newcommand{\vol}{\mathrm{vol}\,}
\title{Cone-Volumes And Semialgebraic Sets}
\author{Tom Baumbach, Martin Henk}
\date{Source: arXiv:2506.15370v2 (CC BY 4.0)}
\begin{document}
\maketitle

\noindent\emph{Source and licence.} Cone-Volumes And Semialgebraic Sets. Version arXiv:2506.15370v2 (CC BY 4.0), distributed by arXiv under the Creative Commons Attribution 4.0 International licence, http://creativecommons.org/licenses/by/4.0/, as recorded in the arXiv licence metadata for this exact version and shown on the abstract page https://arxiv.org/abs/2506.15370v2. Full licence notice: \texttt{licenses/arxiv-2506.15370-CC-BY-4.0.txt}.\par\smallskip

\noindent\textit{Editorial scope.} Counted unit: the article's lemma lem:continuous (source0.tex lines 1135--1141). The article's complete original proof of it is delivered with it (source0.tex lines 1142--1160, one \texttt{proof} environment of 19 lines). No mathematics is written, repaired or re-derived: the statement and the proof are the article's own lines, in the article's own notation, and no retained line is substituted or re-flowed.  No further source-local context is needed: every cross-reference of every retained line resolves inside the two delivered spans.  Nothing was omitted from any retained span, so the package carries no omission record.  No retained line contains a citation, so the wrapper carries no bibliography and no bibliography entry.  No retained line contains a figure environment, an image-inclusion command or drawing code, so no image is delivered and the package declares no asset dependency.  Editorial formatting only, in addition: an AMS \texttt{amsart} preamble that restates the article's own declarations and macros used by the retained lines, namely the environments \texttt{lemma} on separate counters, together with the standard packages those lines need.  The retained environments print in this document as Lemma 1 in order of appearance, whereas the article prints the same items with the numbering of its own full document.  The complete native source of the article as distributed in the arXiv e-print for this exact version is bundled unchanged as \texttt{sources/180079/source0.tex}.
\begin{lemma} Let $b^{(j)},b\in\R^m_{\geq 0}$, $j\in\N$,  with
  $\lim_{j\to\infty} b^{(j)} = b$. Then
  \begin{equation*}
    \lim_{j\to\infty} \gamma(U,b^{(j)})= \gamma(U,b).
  \end{equation*}
\label{lem:continuous}
\end{lemma}

\begin{proof} As $b^{(j)}\to b$  we have $P(U,b^{(j)})\to P(U,b)$ in
  the Hausdorff metric. As $\vol(P(U,b^{(j)}))\to \vol(P(U,b))$ it
  suffices to consider the case $\vol(P(U,b))>0$. Hence, $P(U,b)$ has facets and let $S\subseteq
  U$ be the vectors corresponding to these facets. Moreover, for
  $c\in\R^m$ and $u\in U$ let
  \begin{equation*}
     F(u,c)= P(U,c)\cap\{x\in\R^n : \langle u,x\rangle =c_{\{u\}}\}.
  \end{equation*}
  Then for $u\in S$  and 
 for all sufficiently large $j$, $F(u,b^{(j)})$ is a facet of
 $P(U,b^{(j)})$ and  $F(u,b^{(j)})\to F(u,b)$. Thus 
\begin{equation*}
            \lim_{j\to\infty}  \gamma(U, b^{(j)})_{\{u\}}=  \lim_{j\to\infty} \frac{1}{n} b^{(j)}_{\{u\}}\,\vol_{n-1}(F(u,b^{(j)}))=\gamma(U,b)_{\{u\}}.
\end{equation*}
As $\vol(P(U, b^{(j)})) \to\vol(P(U,b))=\sum_{u\in S} \gamma(U,b)_{\{u\}}$ we conclude for $u\in U\setminus S$
\begin{equation*}
            \lim_{j\to\infty}  \gamma(U, b^{(j)})_{\{u\}}= 0=\gamma(U,b)_{\{u\}}.
\end{equation*}
\end{proof}

\end{document}
